% !TeX root = ../main.tex

\newpage
\section{Fondements Mathématiques : Le Groupe de Möbius}
\label{sec:theorie_mobius}

\subsection{L'isomorphisme entre Géométrie et Algèbre}

La groupes kléinéens se manifestent dans la dualité entre les transformations géométriques du plan complexe et l'algèbre avec des matrices.

Pour étudier formellement les groupes kleinéens sans se heurter aux singularités (divisions par zéro), il est nécessaire d'étendre notre espace de travail géométrique.

\begin{definition}[Sphère de Riemann]
On définit le plan complexe achevé, ou \textit{Sphère de Riemann}, noté $\hat{\mathbb{C}}$, comme la compactification d'Alexandrov du plan complexe, c'est-à-dire l'union de $\mathbb{C}$ et d'un unique point à l'infini :
\[ \hat{\mathbb{C}} = \mathbb{C} \cup \{\infty\} \]
\end{definition}

\begin{remark}[Prolongement et Bijectivité]
Ce cadre topologique est crucial. Pour une transformation de Möbius $f(z) = \frac{az+b}{cz+d}$ (avec $c \neq 0$), l'ajout du point $\infty$ permet de lever la valeur interdite du dénominateur en posant naturellement :
\[ f\left(-\frac{d}{c}\right) = \infty \quad \text{et} \quad f(\infty) = \frac{a}{c} \]
Grâce à ce prolongement, la transformation $f$ devient une bijection parfaite de $\hat{\mathbb{C}}$ dans lui-même. C'est ce qui nous autorise formellement à parler de "groupe" des transformations de Möbius.
\end{remark}\begin{remark}[Interprétation géométrique]
Le terme "compactification d'Alexandrov" signifie topologiquement l'ajout d'un unique point pour rendre un espace compact. Géométriquement, cela s'illustre par la \textbf{projection stéréographique} : on identifie le plan complexe $\mathbb{C}$ à une sphère privée de son pôle Nord. L'ajout du point à l'infini $\infty$ correspond simplement à "reboucher" ce pôle Nord. Ainsi, lorsque $z \to \infty$ dans le plan, son image sur la sphère se rapproche naturellement de ce pôle.
\end{remark}

\begin{definition}[Transformation de Möbius]
On appelle transformation de Möbius toute application $f : \hat{\mathbb{C}} \to \hat{\mathbb{C}}$ de la forme :
\[ f(z) = \frac{az+b}{cz+d} \quad \text{avec} \quad a,b,c,d \in \mathbb{C} \quad \text{et} \quad ad-bc \neq 0 \]
\end{definition}

\begin{theorem}[Représentation matricielle]
L'application $\Phi : SL(2, \mathbb{C}) \to \mathcal{M}$ qui à toute matrice $\begin{pmatrix} a & b \\ c & d \end{pmatrix}$ associe la transformation $z \mapsto \frac{az+b}{cz+d}$ est un morphisme de groupes.
\end{theorem}

\begin{proof}
Soient $M = \begin{pmatrix} a & b \\ c & d \end{pmatrix}$ et $M' = \begin{pmatrix} a' & b' \\ c' & d' \end{pmatrix}$ deux matrices de $SL(2, \mathbb{C})$. 
La composition des transformations associées donne :
\begin{align*}
f_M(f_{M'}(z)) &= \frac{a \left( \frac{a'z+b'}{c'z+d'} \right) + b}{c \left( \frac{a'z+b'}{c'z+d'} \right) + d} \\
&= \frac{(aa'+bc')z + (ab'+bd')}{(ca'+dc')z + (cb'+dd')}
\end{align*}
On reconnaît les coefficients de la matrice produit $M \times M'$. Ainsi, composer deux transformations revient à multiplier leurs matrices représentatives.
\end{proof}

\subsection{Invariants et Trace}
Dans la suite de ce travail, nous utiliserons la \textbf{trace} de ces matrices, notée $\mathrm{Tr}(M) = a+d$, comme invariant fondamental pour déterminer la nature de la transformation (elliptique, parabolique ou hyperbolique) et pour implémenter nos algorithmes d'élagage.